\section{Bài học kinh nghiệm}

Quá trình thực hiện đề tài cho thấy giá trị của việc tách rõ nghiên cứu, thiết kế, triển khai và kiểm thử. Các bài học dưới đây được trình bày ở mức tổng kết, không lặp lại chi tiết kỹ thuật đã phân tích ở các chương trước.

\subsection{Bài học về phân tích bài toán}

\subsubsection{Xác định phạm vi MVP}

Feature Flag là miền bài toán rộng, dễ mở sang A/B testing, realtime sync hoặc workflow phê duyệt. Việc xác định MVP giúp ưu tiên các năng lực cốt lõi: quản lý flag, rule engine xác định, audit log, dashboard và demo tích hợp.

\subsubsection{Phân tách yêu cầu bắt buộc và mở rộng}

Tách yêu cầu bắt buộc khỏi yêu cầu khuyến khích giúp sản phẩm có baseline ổn định trước khi bổ sung SDK, RBAC, statistics, group kill switch và Docker Compose.

\subsubsection{Chuyển yêu cầu thành use case và thiết kế}

Các yêu cầu ban đầu cần được chuyển thành use case, contract API, mô hình dữ liệu và tiêu chí kiểm thử cụ thể; nếu không, hệ thống dễ chỉ dừng ở giao diện bật/tắt đơn giản.

\subsection{Bài học về thiết kế hệ thống}

\subsubsection{Tầm quan trọng của evaluation precedence}

Một rule engine chỉ đáng tin cậy khi thứ tự ưu tiên được định nghĩa trước và được kiểm thử. Đây là điểm quyết định để kết quả evaluation có thể giải thích được.

\subsubsection{Phân tách định nghĩa flag và cấu hình runtime}

Tách metadata của flag khỏi cấu hình theo environment giúp cùng một flag phục vụ nhiều môi trường mà không nhân bản định nghĩa.

\subsubsection{Thiết kế auditability từ đầu}

Audit log nên được đưa vào luồng mutation ngay từ đầu, vì bổ sung sau khi hệ thống đã có nhiều endpoint sẽ dễ bỏ sót hành động quan trọng.

\subsubsection{Cân bằng giữa tính đơn giản và khả năng mở rộng}

Kiến trúc một backend service đủ đơn giản cho mini project nhưng vẫn giữ ranh giới module để có thể tách service hoặc thêm hạ tầng phân tán sau này.

\subsection{Bài học về triển khai và kiểm thử}

\subsubsection{Validation tại ranh giới hệ thống}

DTO validation và constraint database giúp phát hiện lỗi sớm, giảm trạng thái cấu hình không hợp lệ lọt vào engine.

\subsubsection{Kiểm thử các trường hợp precedence}

Các test quan trọng nhất là những test chứng minh precedence, fail-closed và stable rollout không thay đổi khi refactor.

\subsubsection{Quản lý cache invalidation}

Cache chỉ an toàn khi có chiến lược invalidation rõ. Với Feature Flag, snapshot cũ có thể trực tiếp làm sai quyết định runtime.

\subsubsection{Xây dựng dữ liệu và kịch bản demo}

Seed data và kịch bản demo cần được chuẩn bị cùng lúc với tính năng để người xem có thể quan sát giá trị hệ thống mà không phải cấu hình thủ công quá nhiều.
